% s-053-p053-m004; original PDF 53, 54
\Aphorism{s-053-p053-m004}{52}
Wahrheit vergeht, Täuschung besteht. Nur im Selbstbetrug. und Betrug Andrer lässt sich auf die Dauer leben; unverhüllte Zustände können in unserem Klima nur Ausnahmezustände sein. Ich frage, ob z. B. ein einziges Liebes- oder Freundschaftsverhältniss zur Blüthe kommen würde, wenn von Anfang an volle Aufrichtigkeit bestünde. Sobald man zu Zweien oder zu Vielen lebt, ist ein verschliessbarer Schrank und sogar ein geheimes Fach darin unerlässlich. Wir müssen lügen, das heisst soviel als: wir müssen den Schlüssel zu Einigem, das wir sind, in der Hand behalten. Unsere Wahrheitsapostel aus dem Norden, die das Leben auf Offenheit und Vertrauen einrichten wollen: an was für ein unmögliches Leben, an was für unmögliche Menschen denken sie eigentlich dabei? unmöglich, und nicht einmal erstrebenswerth? Denn wie in einem Hause, wo Schlösser und Schlüssel fehlen, Niemand eine Kostbarkeit aufbewahren wird, so ist alles Feine, Seltene, Werthvolle des Lebens darauf angewiesen, dass Lüge und Verheimlichung weiterbesteht. Eine Seele, die dieser Schutzmittel enträth, kann sich keine vornehme Pflege und Bildung angedeihen lassen: auf dem \Foreign{ager publicus} baut man keine edelen und gefährdeten Gewächse an.
\par

